% KYUSHU UNIVERSITY beamer template
% Created by Yu in 2024

\documentclass[11pt]{beamer}
% You can use \documentclass[11pt,aspectratio=169]{beamer} 
% to adjust the aspect ratio to 16:9.

\usepackage{graphicx}
\usepackage{xcolor}
\definecolor{vino}{RGB}{128,0,32}
\usepackage{listings}
\usepackage[
	backend=biber,
	style=verbose,
	sorting=ynt
]{biblatex}

\addbibresource{references.bib}

\usetheme{Madrid}

\author[Grupo 3]{Catalina Leal Rojas, Lucas Daniel Carrillo Aguirre, Lucas Eduardo Veras Costa}
\title[Sección 2]{Comparación entre modelos}

\setbeamertemplate{navigation symbols}{}


% Uncomment the section below to show a table of contents slide every time a new section starts
% \AtBeginSection[]
% {
%  \begin{frame}<beamer>
%  \frametitle{Table of contents}
%  \tableofcontents[currentsection]
%  \end{frame}
% }

\usepackage{Kyushu}

\begin{document}

\begin{frame}
\titlepage
\end{frame}

\section{Example}

\begin{frame}{Nuestro mejor modelo: \textbf{Random Forest}}
\begin{itemize}
\item Utilizamos de base un modelo lineal que incluye las 26 variables de la muestra. Inlcuimos variables que describen las condiciones demográficas, laborales, educativas y de vivienda de los hogares colombianos.
\item Hiperparámetros:
\begin{itemize}
    \item Número de árboles y muestras bootstrap = \textcolor{blue}{1000} (valor fijo)
    \item Número de variables seleccionadas aleatoriamente para cada partición = \textcolor{blue}{6} (de grilla [6,8,9,10] con cross-validation de 5 folds y F1 como métrica para optimizar)
    \item Número mínimo de observaciones por nodo = \textcolor{blue}{1} (valor fijo)
\end{itemize}
\item Desbalance de clases: 
\begin{itemize}
    \item Umbral de clasificación = \textcolor{blue}{0.34} (el valor que maximiza el F1 dentro de muestra entre una grilla (0.05,0.95))
    \item Downsampling
\end{itemize}
\item \textbf{Kaggle score:}\textcolor{blue}{0.63}
\end{itemize}


\end{frame}

\begin{frame}{Top 5 otros modelos}
    \begin{itemize}
        \item Random Forest (Bagging) (\textcolor{blue}{0.63})
        \item Logit (\textcolor{blue}{0.62})
        \item Random Forest (diferente técnica para tratar el desbalance de clases) (\textcolor{blue}{0.61})
        \item Lineal (\textcolor{blue}{0.61})
    \item Elastic Net (\textcolor{blue}{0.60})
    \end{itemize}
\end{frame}

\begin{frame}{Bagging (0.63)}
\begin{itemize}
    \item Hiperparámetros:
\begin{itemize}
    \item Número de árboles y muestras bootstrap = \textcolor{blue}{500} (valor fijo)
    \item Número de variables seleccionadas aleatoriamente para cada partición = \textcolor{blue}{26} (todas las variables)
    \item Número mínimo de observaciones por nodo = \textcolor{blue}{50} (de grilla [1,5,15,50] con cross-validation de 5 folds y F1 como métrica para optimizar)
    \item \textit{Split Rule} = \textcolor{blue}{Hellinger} (de grilla [Gini, Hellinger] con cross-validation de 5 folds y F1 como métrica para optimizar)
\end{itemize}
\item Desbalance de clases: 
\begin{itemize}
    \item Umbral de clasificación = \textcolor{blue}{0.52} (el valor que maximiza el F1 dentro de muestra entre una grilla (0.05,0.95))
    \item Downsampling
\end{itemize}
\end{itemize}
\end{frame}

\begin{frame}{Bagging (0.63)}
    \item \textbf{¿Por qué da el mismo puntaje?:}
    \begin{itemize}
        \item Creemos que se alcanzó una región de estabilidad en el espacio de hiperparámetros (más árboles no mejoran mucho la predicción).
        \item El número de observaciones en cada nodo podría ser un ejemplo del trade-off sesgo-varianza
        \item \textbf{Buena selección de variables:} las variables con mayor poder predictivo tienden a ser seleccionadas de manera recurrente en las primeras divisiones de los árboles (bosques similares, incluso con distintos hiperparámetros). 
        \item El criterio de división Hellinger (pensado para el desbalance de clases) pudo haber compensado la menor complejidad en términos de número de árboles y condiciones sobre los nodos terminales.
    \end{itemize}
\end{frame}

\begin{frame}{Logit (0.62)}
\begin{itemize}
    \item Se utilizó un modelo lineal que incluía las 26 variables creadas por el equipo
    \item Desbalance de clases: 
\begin{itemize}
    \item Umbral de clasificación = \textcolor{blue}{0.21} (el valor que maximiza el AUC (ROC))
\end{itemize}
\end{itemize}
\end{frame}

\begin{frame}{Logit (0.62)}
    \item \textbf{¿Por qué da casi el mismo puntaje?:}
    \begin{itemize}
    \item Variables altamente informativas y bien seleccionadas
    \item Trade-off sesgo-varianza
    \item En ambos se buscaron soluciones para el desbalance de clases. Parece indicar que solamente cambiar el umbral de clasificación puede mejorar sustancialmente cómo el modelo identifica la clase minoritaria ($F_1$)
    \end{itemize}
\end{frame}



\begin{frame}{Random Forest: otra técnica para el desbalance (0.61)}
\begin{itemize}
    \item Mismos hiperparámetros que el modelo ganado
    \item La diferencia radica en que el umbral de clasificación fue elegido con la maximización de la AUC (ROC) igual a \textcolor{blue}{0.21}
    \item En este caso, elegir la misma métrica que evalúa Kaggle, mejora marginalmente el desempeño del modelo.
    
\end{itemize}
    
\end{frame}

\begin{frame}{Modelo de probabilidad lineal (0.61)}

\begin{equation} % usa * si no quieres número
\begin{aligned}
Pobre_i = \beta_0 + \mathbf{X}_i'\boldsymbol{\beta}
+ \beta_1 Nper_i^2 + \beta_2 edad\_head_i^2 \\
+ \beta_3 (bin\_headWoman_i \times cat\_educHead_i) \\
\quad + \beta_4 (edad\_head_i \times cat\_educHead_i)
+ \beta_5 (Nper_i \times t\_dependencia_i)
+ \varepsilon_i
\end{aligned}
\end{equation}

\begin{itemize}
\item Desbalance de clases:
\begin{itemize}
    \item Umbral de clasificación = \textcolor{blue}{0.27} (el valor que maximiza el AUC (ROC))
\end{itemize}

\end{itemize}
\end{frame}

\begin{frame}{Modelo de probabilidad lineal (0.61)}
 
    \begin{itemize}
    \item Para compensar las predicciones fuera de [0,1] y la menor eficiencia del modelo lineal frente a Logit, se incluyeron términos cuadráticos e interacciones para añadir un poco más de complejidad
    \item También creemos que el cálculo de un umbral de clasificación mejor ajustado a predecir los verdaderos positivos logra equiparar un modelo tan sencillo con técnicas como Random Forest
    \end{itemize}
\end{frame}

\begin{frame}{Elastic Net (0.60)}
\begin{itemize}
    \item Hiperparámetros:
\begin{itemize}
    \item $\alpha$ = \textcolor{blue}{0.7} (grilla de 0 a 1, con saltos de 0.1). Este modelo es 70\% Lasso y 30\%  Ridge
    \item $\lambda$ = \textcolor{blue}{0.001} (grilla de $10^{-3}$ a $10^3$, con 10 valores logarítmicamente espaciados)
\end{itemize}
\item Desbalance de clases: 
\begin{itemize}
    \item Umbral de clasificación = \textcolor{blue}{0.2} (el valor que maximiza el F1 dentro de muestra entre una grilla (0.05,0.95))
\end{itemize}
\end{itemize}
\end{frame}

\begin{frame}{Elastic Net (0.60)}
\begin{itemize}
 \item Aunque combina regularización LASSO y Ridge para manejar multicolinealidad y selección de variables, sigue siendo un modelo lineal en los parámetros. Esto, como ya vimos, tiene un desempeño ligeramente menor que otros modelos como Random Forest
 \item Nuevamente, los resultados siguen cercanos a otros modelos más complejos gracias al tratamiento para el desbalance de observaciones entre pobres y no pobres.
\end{itemize}
\end{frame}

\begin{frame}{Los que quedaron por fuera del top}
\begin{itemize}
    \item Gradient Boosting
    \item Elastic net. De este último realizamos una prueba para ver los efectos del remuestreo en este modelo en particular:
    \begin{figure}[H]
\centering
\includegraphics[width=0.6\textwidth]{template/ROC_NaiveBayes_sampling.png}
\label{fig:mi_foto}
\end{figure}

\end{itemize}
\end{frame}

\begin{frame}{Conclusión}
\begin{itemize}
  \item Los modelos presentaron desempeños muy similares, reflejando la alta calidad y relevancia de las variables construidas. 
  \item El \textit{Random Forest} obtuvo el mejor rendimiento (\(F_1 = 0.63\)) por su capacidad para capturar relaciones no lineales e interacciones complejas. 
  \item Modelos más simples como \textit{Logit} y el modelo lineal lograron resultados comparables, mostrando que la selección de variables fue más determinante que la complejidad del algoritmo. 
  \item El tratamiento del desbalance de clases mediante remuestreo y calibración del umbral con \(F_1\) mejoró la identificación de los hogares pobres. 
  \item Una buena preparación de datos y una calibración cuidadosa pueden igualar o superar la complejidad algorítmica en la predicción de pobreza. 
\end{itemize}
\end{frame}

\end{document}
